\section*{Data Lakes y Data Lakehouse}

Un Data Lake y un Data Lakehouse son arquitecturas destinadas a almacenar y
procesar grandes volúmenes de información. Aunque comparten ciertos principios,
se diferencian en el grado de administración, confiabilidad y rendimiento que
ofrecen. El presente trabajo explica sus características, arquitecturas,
ventajas, desafíos y casos de uso principales.

\section*{1. ¿Qué es un Data Lake?}

Un \textit{Data Lake} es un repositorio centralizado y escalable que permite
almacenar grandes cantidades de datos estructurados, semiestructurados y no
estructurados en su formato original o casi original. A diferencia de una base
de datos tradicional, no exige definir toda la estructura de la información
antes de almacenarla.

Generalmente utiliza el principio \textit{schema-on-read}: el esquema se aplica
cuando los datos se consultan o analizan, no cuando se incorporan. Esto permite
conservar información para usos futuros, incluso cuando todavía no se conoce su
propósito analítico.

\section*{2. Características de un Data Lake}

Las principales características de un Data Lake son las siguientes:

\begin{itemize}
    \item Almacena datos estructurados, semiestructurados y no estructurados.
    \item Conserva los datos crudos o con un nivel mínimo de transformación.
    \item Puede manejar grandes volúmenes de información.
    \item Admite cargas por lotes y flujos de datos en tiempo real.
    \item Utiliza almacenamiento distribuido o almacenamiento de objetos en la nube.
    \item Aplica el esquema al momento de la lectura.
    \item Permite utilizar diferentes motores de procesamiento sobre los mismos datos.
    \item Favorece la exploración, la analítica avanzada y el aprendizaje automático.
    \item Puede organizarse en zonas según la calidad, sensibilidad y grado de procesamiento de la información.
    \item Requiere catálogo, metadatos, seguridad y gobierno para evitar que se convierta en un pantano de datos.
\end{itemize}

\section*{3. Arquitectura de un Data Lake}

Una arquitectura típica de Data Lake contiene las siguientes capas:

\begin{enumerate}
    \item \textbf{Fuentes de datos.} Bases de datos, aplicaciones, archivos,
    sensores, redes sociales, registros de sistemas y servicios externos.

    \item \textbf{Ingestión.} Recibe los datos mediante procesos por lotes o
    transmisiones en tiempo real.

    \item \textbf{Almacenamiento.} Conserva la información en sistemas
    distribuidos o en almacenamiento de objetos.

    \item \textbf{Zonas de datos.} La zona de aterrizaje o \textit{raw}
    conserva los datos originales; la zona de trabajo se utiliza para
    experimentación y transformación; la zona sensible protege información
    restringida; y la zona de producción o \textit{gold} contiene datos
    limpios, validados y documentados.

    \item \textbf{Procesamiento.} Incluye los motores que limpian,
    transforman, combinan y analizan los datos.

    \item \textbf{Catálogo y gobierno.} Administra metadatos, permisos,
    linaje, calidad, auditoría y políticas de conservación.

    \item \textbf{Consumo.} Proporciona información a herramientas de
    inteligencia de negocios, reportes, ciencia de datos, inteligencia
    artificial y aplicaciones.
\end{enumerate}

El flujo general puede resumirse como: fuentes, ingestión, almacenamiento por
zonas, procesamiento y consumo. La seguridad, el catálogo y el gobierno son
componentes transversales a todas las capas.

\section*{4. Diferencias entre un Data Lake y un Data Warehouse}

Las diferencias principales son:

\begin{itemize}
    \item \textbf{Estado de los datos:} el Data Lake conserva datos crudos o
    parcialmente procesados, mientras que el Data Warehouse contiene datos
    limpios, estructurados y modelados.

    \item \textbf{Tipos de datos:} el lago admite información estructurada,
    semiestructurada y no estructurada; el almacén trabaja principalmente con
    datos estructurados.

    \item \textbf{Esquema:} el Data Lake suele emplear
    \textit{schema-on-read}; el Data Warehouse utiliza
    \textit{schema-on-write}, por lo que el modelo se define antes de cargar.

    \item \textbf{Procesamiento:} en el lago predominan ELT y las
    transformaciones bajo demanda; en el almacén tradicional se utiliza ETL
    antes de cargar la información.

    \item \textbf{Uso:} el Data Lake es apropiado para exploración, ciencia de
    datos, inteligencia artificial y aprendizaje automático. El Data Warehouse
    está optimizado para reportes, indicadores e inteligencia de negocios.

    \item \textbf{Usuarios:} el lago es utilizado principalmente por
    ingenieros, analistas y científicos de datos; el almacén está orientado a
    analistas de negocio y responsables de la toma de decisiones.

    \item \textbf{Costo y flexibilidad:} el Data Lake ofrece almacenamiento
    económico y flexible. El Data Warehouse normalmente tiene un costo mayor,
    pero ofrece consultas SQL más rápidas y predecibles.
\end{itemize}

\section*{5. Ventajas de los Data Lakes}

\begin{itemize}
    \item Almacenan grandes volúmenes de datos a un costo relativamente bajo.
    \item Admiten prácticamente cualquier tipo y formato de información.
    \item Facilitan la incorporación rápida de datos sin transformaciones previas complejas.
    \item Conservan el detalle original para realizar nuevos análisis.
    \item Favorecen la experimentación y el descubrimiento de patrones.
    \item Son apropiados para inteligencia artificial y aprendizaje automático.
    \item Reducen la pérdida de información que podría ser útil en el futuro.
    \item Permiten utilizar diferentes herramientas sobre los mismos archivos.
    \item Favorecen el autoservicio cuando cuentan con un catálogo y gobierno adecuados.
    \item Pueden escalar de forma elástica en plataformas de nube.
\end{itemize}

\section*{6. Riesgos o desafíos de los Data Lakes}

El principal riesgo es que el repositorio se convierta en un \textit{Data
Swamp} o pantano de datos: un entorno lleno de información oscura, duplicada,
sin documentación y difícil de utilizar.

Otros desafíos importantes son:

\begin{itemize}
    \item Baja calidad o falta de validación de los datos.
    \item Ausencia de metadatos y dificultad para encontrar información.
    \item Acceso indebido a datos sensibles.
    \item Falta de gobierno, linaje y responsables claramente definidos.
    \item Duplicación de conjuntos de datos.
    \item Rendimiento poco predecible para algunas consultas.
    \item Crecimiento descontrolado de los costos de almacenamiento y procesamiento.
    \item Necesidad de personal especializado.
    \item Problemas de soberanía, privacidad y cumplimiento normativo.
    \item Incompatibilidades entre herramientas y formatos.
    \item Falta de participación de los usuarios de negocio.
\end{itemize}

\section*{7. Implementación de un Data Lake}

Una implementación recomendable comprende los siguientes pasos:

\begin{enumerate}
    \item Definir los objetivos del negocio y los casos de uso prioritarios.
    \item Conseguir patrocinio directivo y participación de las áreas involucradas.
    \item Identificar las fuentes y clasificar los datos según su sensibilidad.
    \item Elegir una plataforma local, en la nube o híbrida.
    \item Diseñar las zonas de datos y sus políticas de gobierno.
    \item Crear procesos de ingestión por lotes y en tiempo real.
    \item Implementar catálogo, metadatos, linaje y reglas de calidad.
    \item Configurar seguridad, cifrado, permisos y auditoría.
    \item Proporcionar herramientas de preparación, análisis y autoservicio.
    \item Comenzar con un caso de uso pequeño y medible.
    \item Evaluar los resultados y extender gradualmente la plataforma.
    \item Monitorear costos, calidad, rendimiento y cumplimiento.
\end{enumerate}

La construcción puede resumirse en cuatro etapas generales: levantar la
infraestructura, organizar el lago, habilitar el autoservicio y abrirlo
progresivamente a los usuarios.

\section*{8. ¿Qué es un Data Lakehouse?}

Un \textit{Data Lakehouse} es una arquitectura que combina la flexibilidad,
escalabilidad y bajo costo de un Data Lake con la confiabilidad, el rendimiento
y las capacidades de administración de un Data Warehouse.

El Lakehouse mantiene los datos en almacenamiento económico y, normalmente, en
formatos abiertos, pero incorpora una capa de tablas y metadatos que permite
realizar transacciones, controlar esquemas, gobernar la información y ejecutar
consultas analíticas con mayor rendimiento. De esta manera, una misma plataforma
puede atender inteligencia de negocios, ingeniería de datos, ciencia de datos e
inteligencia artificial.

\section*{9. Principales características de un Data Lakehouse}

\begin{itemize}
    \item Almacenamiento escalable y económico.
    \item Soporte para datos estructurados, semiestructurados y no estructurados.
    \item Uso de formatos abiertos, como Apache Parquet.
    \item Capa de tablas y metadatos sobre los archivos.
    \item Transacciones ACID: atomicidad, consistencia, aislamiento y durabilidad.
    \item Validación y evolución de esquemas.
    \item Historial, versionamiento y \textit{time travel}.
    \item Procesamiento por lotes y en tiempo real.
    \item Consultas SQL de alto rendimiento.
    \item Gobierno, seguridad, auditoría y linaje unificados.
    \item Compatibilidad con BI, ciencia de datos, aprendizaje automático e IA.
    \item Separación entre el almacenamiento y la capacidad de cómputo.
\end{itemize}

Tecnologías como Delta Lake, Apache Iceberg y Apache Hudi proporcionan varias
de estas capacidades. En particular, la capa de metadatos permite controlar qué
archivos forman parte de cada versión de una tabla y ofrecer transacciones ACID,
evolución de esquemas y recuperación de versiones.

\section*{10. Arquitectura de un Data Lakehouse}

Una arquitectura de Data Lakehouse puede organizarse de la siguiente manera:

\begin{enumerate}
    \item \textbf{Fuentes:} aplicaciones, bases de datos, dispositivos,
    archivos y servicios externos.

    \item \textbf{Ingestión:} procesos por lotes, replicación y transmisión en
    tiempo real.

    \item \textbf{Almacenamiento de objetos:} repositorio económico para los
    diferentes tipos de datos.

    \item \textbf{Formato de tablas:} tecnologías como Delta Lake, Apache
    Iceberg o Apache Hudi organizan los archivos en tablas administradas.

    \item \textbf{Capa de metadatos:} controla versiones, transacciones,
    esquemas y archivos.

    \item \textbf{Procesamiento:} motores SQL, Apache Spark, herramientas de
    transmisión y plataformas de aprendizaje automático.

    \item \textbf{Catálogo y gobierno:} administra permisos, linaje,
    auditoría, calidad y descubrimiento.

    \item \textbf{Consumo:} proporciona información para tableros, reportes,
    analítica, IA, aprendizaje automático y aplicaciones.
\end{enumerate}

Con frecuencia se emplea una arquitectura de medallón. La capa bronce conserva
los datos originales; la capa plata contiene datos depurados, integrados y
validados; y la capa oro contiene datos agregados y preparados para el negocio.

\section*{11. Diferencias entre un Data Lake y un Data Lakehouse}

\begin{itemize}
    \item \textbf{Organización:} el Data Lake suele administrar archivos y
    objetos, mientras que el Lakehouse presenta tablas administradas sobre esos
    archivos.

    \item \textbf{Transacciones:} un lago tradicional generalmente no ofrece
    garantías ACID; el Lakehouse sí las proporciona.

    \item \textbf{Esquema y calidad:} en el lago, estos controles dependen de
    procesos adicionales; en el Lakehouse existen mecanismos de aplicación y
    evolución de esquemas.

    \item \textbf{Versionamiento:} el Lakehouse ofrece historial y
    \textit{time travel}, características que no siempre existen en un lago
    tradicional.

    \item \textbf{Rendimiento:} el Data Lake puede presentar un rendimiento
    variable, mientras que el Lakehouse incorpora optimizaciones para consultas
    SQL.

    \item \textbf{Gobierno:} en el lago puede estar fragmentado; el Lakehouse
    busca centralizar catálogo, seguridad, linaje y auditoría.

    \item \textbf{Cargas de trabajo:} el Data Lake es adecuado para
    exploración e IA. El Lakehouse permite utilizar los mismos datos también
    para inteligencia de negocios y reportes empresariales.

    \item \textbf{Complejidad:} un lago básico puede ser más sencillo de
    implementar, mientras que el Lakehouse exige una capa transaccional y de
    gobierno más completa.
\end{itemize}

En síntesis, el Lakehouse no elimina el almacenamiento del lago: agrega
capacidades de administración y confiabilidad directamente sobre él.

\section*{12. Ventajas de los Data Lakehouse}

\begin{itemize}
    \item Unifican inteligencia de negocios, ciencia de datos, procesamiento e IA.
    \item Reducen la necesidad de mantener un Data Lake y un Data Warehouse separados.
    \item Evitan copias innecesarias de información.
    \item Disminuyen procesos ETL y problemas de desactualización.
    \item Mejoran la calidad y consistencia mediante transacciones ACID.
    \item Permiten aplicar y evolucionar esquemas.
    \item Proporcionan historial, auditoría y recuperación de versiones.
    \item Aprovechan almacenamiento económico y escalable.
    \item Ofrecen mejor rendimiento para consultas SQL.
    \item Simplifican la seguridad, el catálogo, el gobierno y el linaje.
    \item Promueven formatos abiertos e interoperabilidad.
    \item Facilitan el procesamiento por lotes y en tiempo real.
\end{itemize}

Al trabajar sobre una única plataforma se reduce la duplicación de datos y la
complejidad operativa causada por mantener un lago y un almacén independientes.

\section*{13. Casos de uso recomendados}

\subsection*{Casos apropiados para un Data Lake}

Se recomienda utilizar un Data Lake cuando:

\begin{itemize}
    \item Se necesita conservar grandes cantidades de datos crudos.
    \item Todavía no se conocen todos los usos futuros de la información.
    \item Existen imágenes, audio, video, registros, documentos o datos de sensores.
    \item Los científicos de datos requieren información con el mayor nivel de detalle.
    \item Se desarrollan proyectos exploratorios de IA o aprendizaje automático.
    \item Se necesita un archivo histórico económico.
    \item El rendimiento SQL transaccional no es el requisito principal.
\end{itemize}

Algunos ejemplos son los datos de Internet de las cosas, registros de
navegación, archivos multimedia, investigación científica, análisis
exploratorio y almacenamiento histórico.

\subsection*{Casos apropiados para un Data Lakehouse}

Se recomienda utilizar un Data Lakehouse cuando:

\begin{itemize}
    \item La organización necesita BI, analítica avanzada e IA sobre los mismos datos.
    \item Los datos deben actualizarse de manera confiable.
    \item Se requieren transacciones, control de esquemas y versionamiento.
    \item Es necesario combinar cargas por lotes y flujos en tiempo real.
    \item Se desea reducir duplicaciones entre el lago y el almacén de datos.
    \item Existen requisitos estrictos de seguridad, auditoría, linaje y calidad.
    \item Diferentes equipos necesitan una fuente de datos común y gobernada.
\end{itemize}

Algunos ejemplos son la detección de fraude, la personalización de clientes, el
mantenimiento predictivo, el análisis financiero, el comercio electrónico, las
telecomunicaciones y el análisis clínico con controles de privacidad.

En conclusión, el Data Lake resulta conveniente para conservar y explorar datos
diversos con máxima flexibilidad. El Data Lakehouse es preferible cuando esa
flexibilidad debe combinarse con confiabilidad, gobierno y rendimiento para
aplicaciones empresariales.
